% Chapter Template

% todo en la parte de datasets se me sugirió mover aquí esto...
% Algunos de los experimentos que no resultaron exitosos comprenden:
% pruebas con solo variables predictivas sin embeddings, demuestran que no es posible capturar eventos como el público en vivo o el habla en los primeros intentos de inferencia de \textit{liveness} y \textit{speechiness}. Las pruebas con embeddings pero sin variables predictivas en general no proporcionan la información de \textit{key}, \textit{tempo}, \textit{loudness}, \textit{mood} o \textit{genres}.
% Pruebas con distintas ventanas resultaron ser ineficientes en el proceso de extracción de más de 60~000 canciones. Pruebas aplicando todos los estadísticos resultaron en un espacio de variables que consume mucho espacio de almacenamiento y con alta correlación.

\chapter{Ensayos y resultados} % Main chapter title

\label{Chapter4} % Change X to a consecutive number; for referencing this chapter elsewhere, use \ref{ChapterX}
Todos los capítulos deben comenzar con un breve párrafo introductorio que indique cuál es el contenido que se encontrará al leerlo.  La redacción sobre el contenido de la memoria debe hacerse en presente y todo lo referido al proyecto en pasado, siempre de modo impersonal.

%todo Dada la complejidad de objetivos como los géneros musicales, 197 clases resultantes, se priorizó el análisis de la precisión \textit{Top-K} para capturar la pluralidad semántica de la música
%----------------------------------------------------------------------------------------
%	SECTION 1
%----------------------------------------------------------------------------------------

\section{Pruebas funcionales del hardware}
\label{sec:pruebasHW}

La idea de esta sección es explicar cómo se hicieron los ensayos, qué resultados se obtuvieron y analizarlos.
